\section{Measurement of Callback Execution Time}
\label{app:exec_time}

The following code below shows a part of the PortAudio callback function for the 
DAW, and how the callback execution is measured using \texttt{std::chrono}:

\lstinputlisting[
    style=cppstyle,
    % linerange={process_delay-process_delay},
    caption={The audio callback function showing how the time is measured},
    label={lst:callback-time}
]{appendices/code/exec_time_code.cpp}


This listing below shows the C++ struct definition of the statistical values. 
This holds the temporary values required to calculate the mean, max standard 
deviation. These variables are encapsulated like this to keep the code clean.


\begin{lstlisting}[style=cppstyle, caption={Audio callback statistics structure}, label={lst:audiocallbackstats}]
struct AudioCallbackStats
{
    std::atomic<double> callbackTimeUs{0.0};
    std::atomic<uint64_t> sumUs{0};
    std::atomic<uint64_t> count{0};
    std::atomic<uint64_t> maxUs{0};

    std::atomic<uint64_t> squareSumUs{0};
};
\end{lstlisting}


Listing C.3. is used to calculate the mean value of the callback time, and Listing C.4 is used to calculate the standard deviation. 
This should be done after PortAudio finishes streaming and the program accumulates the sum of the measured callback times, 
the total number of callback times measured, and the squared values of the callback times.


\begin{lstlisting}[style=cppstyle, caption={Function to calculate the mean of the callback times}, 
    label={lst:mean_callback}]
double AudioEngine::getAvgCallbackTimeUs() const noexcept
{
    return stats.sumUs.load() / (double)stats.count.load();
}
\end{lstlisting}


\begin{lstlisting}[style=cppstyle, caption={Function to calculate the standard deviation of the callback times}, 
    label={lst:stddev_callback}]
double AudioEngine::getStdDevCallbackTimeUs() const noexcept
{
    double n = (double)stats.count.load();

    if (n <= 1.0) // Avoid division by zero or negative variance
        return 0.0;

    double sum = stats.sumUs.load();
    double squareSum = stats.squareSumUs.load();

    // Standard deviation based on the sample formula: sqrt((sum(xi - mu)^2) / N - 1)
    double variance = (squareSum - sum * sum / n) / (n - 1.0);
    return std::sqrt(variance);
}
\end{lstlisting}

And lastly, the following header file shows the timer adapted using RAII method, and it is
designed to work outside the audio engine in a more convenient way \cite{nagi2026timingcpp}.

\lstinputlisting[
    style=cppstyle,
    % linerange={process_delay-process_delay},
    caption={Timer for Measuring Elasped Time using RAII},
    label={lst:raii-timer}
]{../include/misc/Timer.hpp}